\documentclass[10pt,a4paper]{amsart}
\usepackage{amsmath,amssymb,amsthm,graphicx}
% For commutative diagrams we use Xy-pic
\usepackage[all]{xy}

% We use 2cell for 2-commutative diagrams.
\xyoption{2cell}
% Theorem environments.
%
\theoremstyle{plain}
\newtheorem{theorem}[subsection]{Theorem}
\newtheorem{proposition}[subsection]{Proposition}
\newtheorem{lemma}[subsection]{Lemma}

\theoremstyle{definition}
\newtheorem{definition}[subsection]{Definition}
\newtheorem{example}[subsection]{Example}
\newtheorem{exercise}[subsection]{Exercise}
\newtheorem{situation}[subsection]{Situation}

\theoremstyle{remark}
\newtheorem{remark}[subsection]{Remark}
\newtheorem{remarks}[subsection]{Remarks}

\numberwithin{equation}{subsection}

% Macros
%
\def\lim{\mathop{\mathrm{lim}}\nolimits}
\def\colim{\mathop{\mathrm{colim}}\nolimits}
\def\Spec{\mathop{\mathrm{Spec}}}
\def\Hom{\mathop{\mathrm{Hom}}\nolimits}
\def\Ext{\mathop{\mathrm{Ext}}\nolimits}
\def\SheafHom{\mathop{\mathcal{H}\!\mathit{om}}\nolimits}
\def\SheafExt{\mathop{\mathcal{E}\!\mathit{xt}}\nolimits}
\def\Sch{\mathit{Sch}}
\def\Mor{\mathop{\mathrm{Mor}}\nolimits}
\def\Ob{\mathop{\mathrm{Ob}}\nolimits}
\def\Sh{\mathop{\mathit{Sh}}\nolimits}
\def\NL{\mathop{N\!L}\nolimits}
\def\CH{\mathop{\mathrm{CH}}\nolimits}
\def\proetale{{pro\text{-}\acute{e}tale}}
\def\etale{{\acute{e}tale}}
\def\QCoh{\mathit{QCoh}}
\def\Ker{\mathop{\mathrm{Ker}}}
\def\Im{\mathop{\mathrm{Im}}}
\def\Coker{\mathop{\mathrm{Coker}}}
\def\Coim{\mathop{\mathrm{Coim}}}

% Boxtimes
%
\DeclareMathSymbol{\boxtimes}{\mathbin}{AMSa}{"02}

%
% Macros for moduli stacks/spaces
%
\def\QCohstack{\mathcal{QC}\!\mathit{oh}}
\def\Cohstack{\mathcal{C}\!\mathit{oh}}
\def\Spacesstack{\mathcal{S}\!\mathit{paces}}
\def\Quotfunctor{\mathrm{Quot}}
\def\Hilbfunctor{\mathrm{Hilb}}
\def\Curvesstack{\mathcal{C}\!\mathit{urves}}
\def\Polarizedstack{\mathcal{P}\!\mathit{olarized}}
\def\Complexesstack{\mathcal{C}\!\mathit{omplexes}}
% \Pic is the operator that assigns to X its picard group, usage \Pic(X)
% \Picardstack_{X/B} denotes the Picard stack of X over B
% \Picardfunctor_{X/B} denotes the Picard functor of X over B
\def\Pic{\mathop{\mathrm{Pic}}\nolimits}
\def\Picardstack{\mathcal{P}\!\mathit{ic}}
\def\Picardfunctor{\mathrm{Pic}}
\def\Deformationcategory{\mathcal{D}\!\mathit{ef}}
\usepackage{cleveref}
\crefname{originallemma}{original lemma}{original lemmas}
\Crefname{originallemma}{Original lemma}{Original lemmas}
\makeatletter
\newlabel{lemma-constructible-quasi-compact-quasi-separated}{{095E}{source}{}{sourcepointer-1}{}}
\newlabel{lemma-constructible-quasi-compact-quasi-separated@cref}{{[originallemma][0][]095E}{[1][0][]source}{}{}{}}
\makeatother
\title{Constructibility of a sheaf on the small etale site is a Zariski local property}
\author{The Stacks Project Authors}
\date{stacks-project:etale-cohomology.tex:local-snapshot-20261011; GNU FDL 1.2 or later}
\begin{document}
\maketitle
\noindent\textit{Editorial note.} All delivered mathematical text is the author's own, copied from the pinned Stacks Project chapter file etale-cohomology.tex; no statement or proof step was rewritten, reordered or removed.\par\medskip
\noindent\textit{Editorial note.} The proof is the author's complete argument: the reformulation, the restriction to affine opens and the application of the criterion are all inside the counted span.\par\medskip
\noindent\textit{Editorial note.} The label lemma-constructible-local occurs in no other staged Stacks chapter, so no cross-chapter restatement had to be weighed for this card.\par\medskip
\noindent\textit{Editorial note.} The single labelled reference occurring inside the counted spans is carried in assembly.external\_references as a non-load-bearing navigation pointer to the same chapter.\par\medskip
\noindent\textit{Editorial note.} The preamble is the chapter's own native preamble excerpt already used for this source: the Xy-pic setup (native preamble lines 14-18) and the native theorem declarations with the mathematical macros (native preamble lines 154-217).\par\medskip
\noindent\textit{Allowed standard prerequisites (editorial).}\begin{itemize}
\item Constructible sheaf of sets, abelian groups or modules on X\_etale. Same chapter, the definition of constructibility by a partition of X into locally closed constructible strata on which the sheaf is locally constant. Supplies the property shown to be local on X.
\item Restriction of a constructible sheaf to an open subscheme. Same chapter, the immediate fact that the restriction of a constructible sheaf to an open subscheme is constructible. Supplies the constructibility of the restrictions to the members of a refinement of the given Zariski covering.
\item Zariski local nature of constructibility on a quasi-compact quasi-separated scheme. Same chapter, the criterion producing a partition of an affine scheme from constructible restrictions on the members of a finite open covering. Supplies the final step, which assembles the local partitions into a global one.
\end{itemize}\medskip
\ifdefined\hypertarget\hypertarget{sourcepointer-1}{}\fi
\noindent\textit{Source reference.} Original lemma 095E: sources/stacks-etale-cohomology/pinned-original-source.tex, source lines 11717--11717. Same-chapter navigation pointer invoked at proof lines 11801-11802, where the author uses the criterion for constructibility on a quasi-compact quasi-separated scheme in terms of finitely many strata. The pointer only resolves that reference for the reader; the cited lemma is not reproduced inside the counted span, and no bibliography citation is involved.\par\medskip
\begin{lemma}
\label{lemma-constructible-local}
Let $X$ be a scheme. Checking constructibility of a sheaf
of sets, abelian groups, $\Lambda$-modules (with $\Lambda$ Noetherian)
can be done Zariski locally on $X$.
\end{lemma}

\begin{proof}
The statement means if $X = \bigcup U_i$ is an open covering
such that $\mathcal{F}|_{U_i}$ is constructible, then $\mathcal{F}$
is constructible. If $U \subset X$ is affine open, then
$U = \bigcup U \cap U_i$ and $\mathcal{F}|_{U \cap U_i}$ is constructible
(it is trivial that the restriction of a constructible sheaf to
an open is constructible). It follows from
Lemma \ref{lemma-constructible-quasi-compact-quasi-separated}
that $\mathcal{F}|_U$ is constructible, i.e., a suitable partition
of $U$ exists.
\end{proof}

\end{document}
